%%%%%%%%%%  scan idx 464  |  folio 451  %%%%%%%%%%
\[
  \lambda^p(J',u')(1-J')=\lambda^p(u'(1-J')).
\]
Comme \(1-J'=\lambda_{-1}(J')\), c'est la formule qui sert à définir \(\lambda^p(J',u')\).

ii) On se place dans le cas général. On compose alors \(i\) avec une section : \(X\to\mathbf P_X^1\), puis avec une section : \(\mathbf P_X^1\to\mathbf P_{\mathbf P_X^1}^1\), de façon que l'immersion composée vérifie la condition de 4.1. (voir 4.2). La proposition 4.3 est alors vraie pour l'immersion composée d'après ce qui précède; il reste donc à « redescendre » à \(X\). On redescend d'abord à \(\mathbf P_X^1\), puis à \(X\), de sorte que la démonstration se fait en appliquant deux fois le lemme suivant (dû à J.P. Jouanolou) :

\textbf{Lemme} : Dans les conditions de 4.3, soit \(s:X\to\mathbf P_X^1\) la section définie en 4.2. Si la proposition 4.3 est vraie pour \(s\circ i\), elle est vraie pour \(i\).

Si on suppose 4.3 vraie pour \(s\circ i\), on a pour tout \(p\ge1\), et tout \(y\in K^\circ(Y)\) :
\begin{equation}\tag{4.3.3}
 (s\circ i)_*\bigl(\lambda^p(N+1,y)\bigr)
 \equiv
 \lambda^p\bigl((s\circ i)_*(y)\bigr)
 \quad\text{mod.\ torsion}.
\end{equation}

La section \(s\) est définie par un idéal isomorphe à \(\mathcal O_p(-1)\) (4.2) ; on note \(\xi\) sa classe dans \(K^\circ(\mathbf P_X^1)\). On a alors la relation suivante pour tout \(x\in K^\circ(X)\) :
\[
  s_*(x)=p^*(x)(1-\xi),
\]
\(p\) étant le morphisme structural \(\mathbf P_X^1\to X\). En effet, comme \(p\circ s=\Id_X\), on a \(x=s^*(p^*(x))\), d'où \(s_*(x)=s_*(s^*(p^*(x)))=p^*(x)s_*(1)\); comme \(s_*(1)=1-\xi\), on a la relation voulue :

La relation (4.3.3) s'écrit alors :
\[
\begin{aligned}
(s\circ i)_*\bigl(\lambda^p(N+1,y)\bigr)
&\equiv \lambda^p\bigl(p^*(i_*(y))(1-\xi)\bigr)\quad\text{mod.\ torsion}\\
&\equiv \lambda^p(\xi,p^*(i_*(y)))(1-\xi)\quad\text{mod.\ torsion}.
\end{aligned}
\]
